\documentclass[aspectratio=169,11pt]{beamer}
\usepackage[utf8]{inputenc}
\usepackage[T1]{fontenc}
\usepackage[english]{babel}
\usepackage{amsmath,amssymb}
\usepackage{graphicx}
\usepackage{url}
\usetheme{default}
\usecolortheme{seahorse}
\setbeamertemplate{navigation symbols}{}
\setbeamertemplate{footline}[frame number]
\setbeamerfont{frametitle}{size=\large}

\title{AI, Human Cognition and Knowledge Collapse}
\subtitle{Acemoglu, Kong \& Ozdaglar (2026)}
\author{William Gradost}
\date{Artificial Intelligence and Economic Modeling \textperiodcentered{} Repository 4}

\begin{document}

\begin{frame}[plain]
  \titlepage
  \begin{center}\small
    \url{https://github.com/wbgradost/ai-04-acemoglu}
  \end{center}
\end{frame}

\begin{frame}{1 \textperiodcentered{} The paper and the agent's problem}
  \textbf{Question.} Can accurate agentic AI improve today's decisions while
  eroding the human learning that sustains tomorrow's collective knowledge?

  \smallskip
  A short-lived, atomistic agent chooses \(e\geq0\), taking public precision
  \(X_t\) and agentic precision \(\tau_A\) as given. Effort produces both
  context-specific and ``thin'' public information, but the latter is an
  externality.
  \[
    Y=\sigma^{-2}+\lambda_Ie+\tau_A \tag{4}
  \]
  Under Assumption 1 (\(\Delta_I=0,\ \Delta_X>0\)),
  \[
  \max_{e\geq0}\;
  f(0,0)+G(X_t)\Delta_G+G(X_t)G(Y)\Delta_X
  -\frac{\varepsilon}{\varepsilon+1}e^{(\varepsilon+1)/\varepsilon}. \tag{6}
  \]
  \footnotesize\(G(\tau)=2\Phi(\sqrt\tau)-1\): probability that a Gaussian
  prediction error with precision \(\tau\) lies within the success band.
\end{frame}

\begin{frame}{2 \textperiodcentered{} Main result and its conditions}
  \textbf{Observation 1 (pp. 14--15).} For positive precisions, \(\lambda_I>0\),
  and Assumption 1,
  \[
  U_e=\Delta_XG(X_t)\lambda_Ig(Y)-e^{1/\varepsilon},
  \]
  \[
  U_{eX}=\Delta_X\lambda_Ig(X_t)g(Y)>0,\qquad
  U_{e\tau_A}=\Delta_XG(X_t)\lambda_Ig'(Y)<0,
  \]
  since \(g=G'>0\) and
  \(g'(\tau)=-\tfrac12(1+\tau^{-1})g(\tau)<0\).

  \smallskip
  \textbf{Mechanism.} General knowledge makes context-specific learning useful;
  AI supplies context-specific precision directly, where returns are diminishing.

  \smallskip
  \textbf{Dynamic condition.} With \(\varepsilon>4\), zero knowledge is locally
  stable; if \(\tau_A>\tau_A^c\), it is uniquely and globally stable
  (Proposition 5). If \(\varepsilon<4\), every \(X_1>0\) converges instead to
  the high-knowledge state (Proposition 3).
\end{frame}

\begin{frame}{3 \textperiodcentered{} What I did}
  \begin{itemize}\small
    \item Reconstructed equations (3)--(6) and the effort FOC from the May 5,
      2026 MIT copy (69 pages), keeping equation and page references.
    \item Derived both cross-partials from
      \(g(\tau)=(2\pi)^{-1/2}e^{-\tau/2}\tau^{-1/2}\), including the corner
      \(X_t=0\Rightarrow e=0\).
    \item Traced the feedback loop:
      \[
      \tau_A\uparrow\Rightarrow e_t\downarrow\Rightarrow
      X_{t+1}\downarrow\Rightarrow e_{t+1}\downarrow.
      \]
    \item Audited welfare: under Assumption 2,
      \(\sigma^{-2}\geq\sqrt2-1\), Propositions 10--11 imply a finite,
      single-peaked welfare optimum in AI precision.
    \item Checked Section 5 and the appendix before proposing an extension:
      aggregation, synthetic data, and effort separability are already covered;
      the production restriction \(\Delta_I=0\) is not relaxed.
  \end{itemize}
\end{frame}

\begin{frame}{4 \textperiodcentered{} Where I did not believe the AI}
  \begin{columns}[T,onlytextwidth]
    \begin{column}{0.54\textwidth}
      \small
      \textbf{First-pass intuition to audit:} ``More accurate information must
      raise welfare.''

      \medskip
      \textbf{Check.} The handwritten derivation will verify the FOC,
      \(g'(Y)<0\), and the direct/indirect welfare decomposition:
      \begin{align*}
      \frac{\partial\bar U^+}{\partial\tau_A}&=DE+IE,\\
      DE&=g(\bar Y_h)G(\bar X_h)\Delta_X\geq0,\\
      IE&=\frac{\partial G(\bar X_h)}{\partial\tau_A}
      [\Delta_G+G(\bar Y_h)\Delta_X]\leq0.
      \end{align*}

      \textbf{Verdict: false dynamically.} Under Assumption 2, welfare rises
      only up to \(\tau_A^*\), then falls (Props. 10--11, p. 27).
    \end{column}
    \begin{column}{0.42\textwidth}
      \centering
      \IfFileExists{hand/observation-1-foc.jpg}{%
        \includegraphics[width=\linewidth,height=.67\textheight,keepaspectratio]
        {hand/observation-1-foc.jpg}%
      }{%
        \fbox{\parbox[c][.55\textheight][c]{.9\linewidth}{\centering
        Handwritten photo pending\\[1ex]
        \texttt{hand/observation-1-foc.jpg}}}%
      }
    \end{column}
  \end{columns}
\end{frame}

\end{document}
